\chapter{Worked example: secret sharing}

Following Rosulek, \emph{The Joy of Cryptography} (secret sharing). In the
2-out-of-2 additive scheme a secret $s$ is split into shares $(k, k \oplus s)$ for
uniform $k$. Either share alone is uniformly distributed and independent of the
secret, so a party holding one share learns nothing --- perfect privacy, the
one-time-pad argument again.

\begin{definition}[Additive 2-of-2 sharing]
  \label{def:ss-share}
  \uses{def:spcomp}
  \lean{CatCrypt.Examples.SecretSharing.revealShare}
  \leanok
  The secret $s$ is shared as $(k, k \oplus s)$ for a uniform key $k$; party $0$
  holds $k$ and party $1$ holds $k \oplus s$.
\end{definition}

\begin{lemma}[Reconstruction]
  \label{lem:ss-reconstruct}
  \uses{def:ss-share}
  \lean{CatCrypt.Examples.SecretSharing.shareXor_reconstruct}
  \leanok
  The two shares reconstruct the secret: $k \oplus (k \oplus s) = s$.
\end{lemma}

\begin{theorem}[Perfect privacy]
  \label{thm:ss-privacy}
  \uses{def:ss-share}
  \lean{CatCrypt.Examples.SecretSharing.ss_perfect_privacy}
  \leanok
  For either corrupted party, any pair of secrets, and any adversary, the privacy
  advantage from a single share is exactly $0$: one share reveals nothing about
  the secret.
\end{theorem}
